\Needspace{8\baselineskip}
\Article{289}
From the general theory given in Article 282, we obtain the following
results on proper representations of binary forms by the ternary form $xx+yy+zz=f$.

I.\ A binary form $\varphi$ cannot be properly represented by $f$
unless it is a primitive positive form and $-1$ (that is, the determinant
of the form $f$) is its characteristic number. So, for a positive
determinant, and also for a negative determinant $-M$ when $M$ is
either divisible by $4$ or of the form $8n+7$, there are no binary
forms properly representable by $f$.

II.\ But if $\varphi=(p,q,r)$ is a primitive positive form of determinant $-M$,
and $-1$ is a characteristic number of the form $\varphi$, and therefore
also of its opposite $(p,-q,r)$, there will be proper representations of
the form $\varphi$ by $f$ belonging to any given value of the expression
$\surd{}-(p,-q,r)$. For all coefficients of the ternary form $g$ of determinant
$-1$ (Article 283) will necessarily be integers, and $g$ will be
a definite form, so it is certainly equivalent to $f$ (Article 285, I).

III.\ By Article 283, III, the number of all representations belonging
to the same value of the expression $\surd{}-(p,-q,r)$ is, in every case except
$M=1$ and $M=2$, as large as the number of transformations of the
form $f$ into $g$, and so, by Article 285, $=48$. The same
article shows that, given one representation belonging to a given value,
we derive the other $47$ from it by permuting the values of $x$,
$y$, $z$
\SourceLeaf{346}
among themselves and reversing their signs in every possible way.
So all $48$ representations give \textit{a single} decomposition of
the form $\varphi$ into three squares, if we consider only the squares
themselves, without regard to their order or the signs of their roots.

IV.\ Let the number of all distinct odd primes dividing $M$ be
$=\mu$. It follows without difficulty from Article 233 that the number
of all distinct values of the expression $\surd{}-(p,\allowbreak -q,\allowbreak r)(\mathrm{mod.}\ M)$ will be $=2^\mu$. By
Article 283, we need consider only half of these (when $M>2$).
So the number of all proper representations of the form $\varphi$ by
$f$ will be $=48\,\allowbreak  .\,\allowbreak  2^{\mu-1}=3\,\allowbreak  .\,\allowbreak  2^{\mu+3}$, while the number of distinct decompositions
into three squares will be $=2^{\mu-1}$.

\textit{Example.} Let $\varphi=19tt+6tu+41uu$, so $M=770$. Here we must consider the
following four values of the expression $\surd{}-(19,\allowbreak -3,\allowbreak 41)(\mathrm{mod.}\ 770)$ (Article 283): $(39,237)$,
$(171,-27)$, $(269,-83)$, $(291,-127)$. To find the representations belonging to the
value $(39,237)$, we first obtain the ternary form $\left(\begin{smallmatrix}19,&41,&2\\3,&6,&3\end{smallmatrix}\right)=g$. By the rules
of Articles 272, 275, we find that $f$ goes into it by the substitution
$\left\{\begin{smallmatrix}1,&-6,&-0\\-3,&-2,&-1\\-3,&-1,&-1\end{smallmatrix}\right\}$, which gives the following representation of the form $\varphi$
by $f$:
\[
x=t-6u,\quad y=-3t-2u.\quad z=-3t-u
\]
For brevity, we do not write out the other $47$ representations
belonging to the same value, which arise by permuting these values and
reversing their signs. All $48$ representations give the same
decomposition of the form $\varphi$ into the three squares
\[
tt-12tu+36uu,\quad 9tt+12tu+4uu,\quad 9tt+6tu+uu
\]

In exactly the same way, the value $(171,-27)$ gives the decomposition into
the squares $(3t+5u)^2$, $(3t-4u)^2$, $tt$; the value $(269,-83)$ gives $(t+6u)^2+(3t+u)^2+(3t-2u)^2$;
and finally the value $(291,-127)$ gives $(t+3u)^2+(3t+4u)^2+(3t-4u)^2$. Each of these decompositions
is equivalent to $48$ representations. --- Besides these $192$
representations, or four decompositions, there will be no others, since
$770$ is divisible by no square, so improper representations cannot exist.

\SourceLeaf{347}
\Needspace{8\baselineskip}
\Article{290}
We will briefly discuss separately the forms of determinants $-1$
and $-2$, which were subject to certain exceptions. We first make
the general observation that, if $\varphi$, $\varphi'$ are any equivalent
binary forms, and $(\Theta)$ is a given transformation of the former into
the latter, combining any representation of the form $\varphi$ by some
ternary form $f$ with the substitution $(\Theta)$ gives a representation
of the form $\varphi'$ by $f$. Further, proper representations of
$\varphi$ give proper representations of $\varphi'$ in this way; distinct
ones give distinct ones, and all of them give all of them. All this is
very easily checked by calculation. So either of the forms $\varphi$,
$\varphi'$ can be represented by $f$ in as many ways as the other.

I.\ First let $\varphi=tt+uu$, and let $\varphi'$ be any other positive binary
form of determinant $-1$, to which $\varphi$ will therefore be
equivalent. Let $\varphi$ go into $\varphi'$ by the substitution $t=\alpha t'+\beta u'$,
$u=\gamma t'+\delta u'$. The form $\varphi$ is represented by the ternary form $f=zz+yy+zz$
by setting $x=t$, $y=u$, $z=0$. Permuting $x,y,z$ gives
\textit{six} representations, and changing the signs of $t$,
$u$ gives four from each, making $24$ distinct representations
in all. A single decomposition into three squares is equivalent to
these, and it is easy to see that there can be no other representations.
It follows that the form $\varphi'$, too, can be decomposed into three
squares in only one way, namely $(\alpha t'+\beta u')^2$, $(\gamma t'+\delta u')^2$ and $0$.
This decomposition is equivalent to $24$ representations.

II.\ Let $\varphi=tt+2uu$, and let $\varphi'$ be any other positive binary
form of determinant $-2$, into which $\varphi$ goes by the substitution
$t=\alpha t'+\beta u'$, $u=\gamma t'+\delta u'$. We then conclude, as in the preceding case, that
$\varphi$, and therefore also $\varphi'$, can be decomposed into three
squares in only one way: $\varphi$ into $tt+uu+uu$, and $\varphi'$ into
$(\alpha t'+\beta u')^2+(\gamma t'+\delta u')^2+(\gamma t'+\delta u')^2$. It is easy to see that such a decomposition is equivalent
to $24$ representations.

It follows that binary forms of determinants $-1$ and $-2$
agree completely with other binary forms as to the number of their
representations by the ternary form $xx+yy+zz$. For in both cases
$\mu=0$, and the formula given in IV of the preceding article does
indeed give $24$ representations. The reason is that the two
exceptions affecting these forms compensate for each other.

For brevity, we omit the application to the form $xx+yy+zz$ of the general
theory of improper representations developed in Article 284.

\SourceLeaf{348}
\Needspace{8\baselineskip}
\Article{291}
By Article 281, the problem of finding all proper representations of
a given positive \textit{number} $M$ by the form $xx+yy+zz$ first
reduces to finding the proper representations of the number $-M$
by the form $-xx-yy-zz=f$. By the rules of Article 280, these are obtained
as follows.

I.\ Find all classes of binary forms of determinant $-M$ whose
forms can be properly represented by $XX+YY+ZZ=F$ (the ternary form whose
adjoint is $f$). When $M\equiv0,\allowbreak \ 4\text{ or }7\ (\mathrm{mod.}\ 8)$, Article 288 shows that there are
no such classes, so $M$ cannot be decomposed into three squares
having no common divisor\HistoricalNote{*)}{This impossibility also
follows from the fact that the sum of three odd squares is necessarily
$\equiv3\ (\mathrm{mod.}\ 8)$; the sum of two odd squares and one even square is either
$\equiv2$ or $\equiv6$; the sum of one odd square and two even squares
is either $\equiv1$ or $\equiv5$; and finally the sum of three even
squares is either $\equiv0$ or $\equiv4$. But in the last case the
representation is clearly improper.}. But when $M\equiv1,\ 2,\ 5\text{ or }6$, there will
be a positive properly primitive genus, and when $M\equiv3$, an improperly
primitive genus, containing all those classes. Let $k$ denote
the number of these classes.

II.\ Now choose $k$ forms at will from these classes, one from
each, and call them $\varphi$, $\varphi'$, $\varphi''$, and so on. Find all
proper representations of all of them by $F$. Their number will
therefore be $3\, .\, 2^{\mu+3}k=K$, where $\mu$ denotes the number of (odd)
prime factors of $M$. Finally, from any such representation, as
\[
X=mt+nu,\quad Y=m't+n'u,\quad Z=m''t+n''u
\]
derive the following representation of $M$ by $xx+yy+zz$:
\[
x=m'n''-m''n',\quad y=m''n-mn'',\quad z=mn'-m'n
\]
The collection of these $K$ representations, which we denote by
$\Omega$, will necessarily contain all representations of $M$.

III.\ So we only need to ask whether \textit{identical} representations
can occur in $\Omega$. Article 280, III already shows that representations
in $\Omega$ derived from different forms, \textit{for example} from
$\varphi$ and $\varphi'$, must be
\SourceLeaf{349}
distinct. So the only remaining question is whether distinct representations
of the same form, \textit{for example} $\varphi$, by $F$ can give
identical representations of the number $M$ by $xx+yy+zz$. It is
immediately clear that, if the representations of $\varphi$ include
\[
X=mt+nu,\quad Y=m't+n'u,\quad Z=m''t+n''u\quad\ldots\ (r)
\]
they also include
\[
X=-mt-nu,\quad Y=-m't-n'u,\quad Z=-m''t-n''u\quad\ldots\ (r')
\]
and that both give the same representation of $M$, which we
denote by $(R)$. So let us examine whether this same $(R)$ can
also follow from other representations of the form $\varphi$. From
Article 280, III, on setting $\chi=\varphi$ there, we readily deduce that,
if all proper transformations of the form $\varphi$ into itself are
given by
\[
t=\alpha t+\beta u,\quad u=\gamma t+\delta u
\]
then all representations of the form $\varphi$ from which $R$
follows will be expressed by
\[
\begin{aligned}
x&=(\alpha m+\gamma n)t+(\beta m+\delta n)u\\
y&=(\alpha m'+\gamma n')t+(\beta m'+\delta n')u\\
z&=(\alpha m''+\gamma n'')t+(\beta m''+\delta n'')u
\end{aligned}
\]
But the theory of transformations of binary forms of negative determinant
developed in Article 179 shows that, in every case except $M=1$
and $M=3$, there are only two proper transformations of the form
$\varphi$ into itself: $\alpha,\allowbreak \beta,\allowbreak \gamma,\allowbreak \delta=1,\allowbreak 0,\allowbreak 0,\allowbreak 1$ and $=-1,0,0,-1$, respectively. (Since
$\varphi$ is a primitive form, what was denoted by $m$ in Article
179 will be either $1$ or $2$, so, apart from the exceptional
cases, $(1)$ certainly applies there.) Thus $(R)$ can arise
only from $r$, $r'$, so every proper representation of the
number $M$ occurs twice, and no more, in $\Omega$. The number
of all distinct proper representations of $M$ will therefore
be $\tfrac{1}{2}K=3\,\allowbreak  .\,\allowbreak  2^{\mu+2}k$.

As for the exceptional cases, Article 179 gives the number of proper
transformations of the form $\varphi$ into itself as $4$ for $M=1$, and
$6$ for $M=3$. And it is easily confirmed that the numbers
of proper representations of the numbers $1$, $3$ are
$\tfrac{1}{4}K$, $\tfrac{1}{6}K$, respectively. For each number can be decomposed
into three
\SourceLeaf{350}
squares in only one way: $1$ into $1+0+0$, and $3$ into
$1+1+1$. The decomposition of $1$ gives six distinct representations,
and that of $3$ gives eight. But $K$ is $=24$ for $M=1$
(where $\mu=0$, $k=1$), and $=48$ for $M=3$ (where $\mu=1$,
$k=1$).

Further, let $h$ denote the number of classes in the principal
genus, which by Article 252 equals the number of classes in any other
properly primitive genus. We observe that $k=h$ for $M\equiv1,\allowbreak \ 2,\allowbreak \ 5\text{ or }6\ (\mathrm{mod.}\ 8)$, but
$k=\tfrac{1}{3}h$ for $M\equiv3\ (\mathrm{mod.}\ 8)$, with the single exception of $M=3$, where
$k=h=1$. So, for numbers of the form $8n+3$, the number of representations
is \textit{in general} $=2^{\mu+2}h$, since for the number $3$ the two
exceptions compensate for each other.
